\section{Model Formulation}
\label{sec:model_formulation}

\placeholder{本研究ではV2G契約が個人の活動スケジューリングに与える制約を評価するための一貫したフレームワークを提供する。}
\placeholder{\citet{wang2025joint}の活動スケジューリングモデルを拡張し、V2Gサービスを考慮した制約条件を導入する。}
\placeholder{V2Gが制約する効果を捉えることができる}
\placeholder{我々は、..., xxについて説明する}

\subsection{OASIS型\citep{pougala2023oasis}の定式化}
\placeholder{OASIS型の定式化を説明する。duwn/duskの定義、活動の定義、効用関数の定義、制約条件の定義を説明する。}

\subsection{活動の不確実性を考慮したwang et al. の拡張}
\subsubsection{\citet{wang2025joint}による活動スケジューリングフレームワーク}
\placeholder{効用関数を説明する}

\subsubsection{活動の発生確率}
\placeholder{活動の発生確率がベルヌーイ分布に従う}

\subsubsection{効用関数}
\placeholder{計画/実行の2段階に分ける}
\placeholder{計画段階は期待効用の最大化、実行段階は効用の最大化を目的とする。}

\subsection{決定変数}
\label{subsec:decision_variable}
\placeholder{決定変数を説明する。}

\subsection{V2Gサービスを考慮した拡張済み制約条件}
\label{subsec:constraints_of_the_v2g_contract}

\subsubsection{活動の連続・順序・時間制約}
\placeholder{活動の連続・順序・時間制約を説明する。}

\subsubsection{SOC, 充電, 放電制約}
\placeholder{SOC, 充電, 放電制約を説明する。}

\subsubsection{V2G契約条件}
\placeholder{V2G契約条件を説明する。}

\subsection{最適化フレームワーク}
\placeholder{離散変数を固定したKKT条件→shadow price, option valueの評価}